\section*{Bab \ref{ch:g12:past-climates} --- Membaca Iklim Masa Lalu}

\begin{solution}{exo:g12:past-climates:1}
(a) \omterm{def:g12:past-climates:proxy}{proksi}; (b) \omterm{def:g12:past-climates:ice-core}{inti es}; (c) \omterm{def:g12:past-climates:glacial-cycle}{interglasial}; (d) \omterm{def:g12:past-climates:pollen}{analisis polen};
(e) \omterm{def:g12:past-climates:lgm}{Maksimum Glasial Terakhir}.
\end{solution}

\begin{solution}{exo:g12:past-climates:2}
(a) Es kutub: karbon dioksida dan metana udara pada masa itu. (b) Sedimen
danau atau gambut: vegetasi, di sini pohon-pohon yang memerlukan musim panas
yang hangat. (c) Sedimen samudra: volume es di daratan dan suhu \omterm{def:g10:layered-ocean:thermocline}{samudra dalam}.
(d) Pohon: kondisi tumbuh, dari tahun ke tahun.
\end{solution}

\begin{solution}{exo:g12:past-climates:3}
(a) $-50 - (-120) = \qty{70}{m}$ dalam \num{7000} tahun: $\qty{70000}{mm}/\num{7000}
= \qty{10}{mm}$ setahun. (b) $10/3.7 \approx 2.7$: kenaikan rata-rata selama
deglasiasi sekitar 2.7 kali lebih cepat daripada kenaikan pada 2006--2018.
\end{solution}

\begin{solution}{exo:g12:past-climates:4}
Satu abad adalah sepuluh dekade: $\qty{2.0}{\degreeCelsius}$ per abad. Seribu
tahun adalah seratus dekade: $\qty{20}{\degreeCelsius}$ per seribu tahun.
\end{solution}

\begin{solution}{exo:g12:past-climates:5}
(a) \num{23000} tahun: presesi sumbu; \num{42000} tahun: kemiringan sumbu;
sekitar \num{100000} tahun: bentuk orbit.
(b) $\num{1000000}/\num{41000} \approx 24$ siklus; $\num{1000000}/\num{100000}
= 10$ siklus.
\end{solution}

\begin{solution}{exo:g12:past-climates:6}
(a) $266.6 + 7.83 \times (-9.4) \approx 266.6 - 73.6 = \qty{193}{ppm}$: hanya
\qty{3}{ppm} di atas \qty{190}{ppm} yang terukur. (b) $266.6 + 7.83 \times 5.5
\approx \qty{310}{ppm}$, \qty{11}{ppm} di atas nilai terukur tertinggi: titik
terhangat terletak di bawah garis. (c) $\Delta T = (427 - 266.6)/7.83 \approx
\qty{20}{\degreeCelsius}$. Garis itu hanya menggambarkan bagaimana kedua besaran
itu berubah bersama, dalam siklus-siklus alami, di dalam rentang datanya
(\qty{-11}{\degreeCelsius} sampai \qty{+6}{\degreeCelsius}, 174 sampai
\qty{299}{ppm}); garis itu bukan hukum, dan nilai yang jauh di luar rentang itu
tidak dapat dibaca darinya.
\end{solution}

\begin{solution}{exo:g12:past-climates:7}
(a) Herba, \qty{84}{\percent}: bentang alam terbuka yang hampir tanpa pohon,
dingin. (b) Ek (\qty{42}{\percent}) dan hazel (\qty{34}{\percent}): hutan daun
lebar, dengan musim panas yang hangat. (c) Birch menyebar lebih dahulu, lalu
pinus; hazel dan ek datang paling akhir. (d) Pohon hanya menyebar secepat
bijinya berpindah dan pohon mudanya tumbuh, dari tempat-tempat ia bertahan
selama masa dingin; tanah juga harus terbentuk. Hutan dapat tertinggal di
belakang iklim.
\end{solution}

\begin{solution}{exo:g12:past-climates:8}
(a) \qty{300}{cm}, yaitu \qty{3000}{mm}, dalam \num{8000} tahun: $3000/8000 =
\qty{0.375}{mm}$ setahun. (b) $\num{6000} + (470 - 300)/300 \times \num{8000}
\approx \num{6000} + \num{4530} \approx \num{10500}$ tahun. (c) $10/0.375
\approx 27$ tahun.
\end{solution}

\begin{solution}{exo:g12:past-climates:9}
(a) $7.16 - 4.11 = 3.05$ juta kilometer persegi, dan $3.05/4.11 \approx
\qty{74}{\percent}$ daratan lebih banyak. (b) Ya: pada \omterm{def:g12:past-climates:lgm}{Maksimum Glasial
Terakhir} permukaan laut berada 125 sampai \qty{134}{m} di bawah permukaan
sekarang, jauh lebih dari \qty{36}{m} yang menyatukan Kalimantan dengan Asia
pada peta itu. (c) Peta itu memakai dasar laut sekarang, yang sejak itu tidak
tetap persis sama (sedimen telah diendapkan di atasnya, dan tanahnya sedikit
bergerak); kisinya juga kasar, satu titik setiap 0.1 derajat, dan aras laut
terendah itu hanya diketahui dengan ketelitian sekitar sepuluh meter.
\end{solution}

\begin{solution}{exo:g12:past-climates:10}
(a) Nilai \omterm{def:g12:past-climates:lgm}{Maksimum Glasial Terakhir}, \qty{4.92}{\permil}, yang lebih tinggi.
Hal ini memang diharapkan: lapisan es yang besar pada masa itu menyimpan air
yang miskin oksigen-18, sehingga samudra, dan cangkang-cangkang yang tumbuh di
dalamnya, menjadi lebih kaya oksigen-18; air dalam juga lebih dingin. (b) Agar
masa-masa yang dingin dan penuh es terplot di bawah dan masa-masa hangat di
atas, seperti pada kurva suhu.
\end{solution}

\begin{solution}{exo:g12:past-climates:11}
(a) Korelasi menunjukkan bahwa kedua rekaman itu naik dan turun bersama, dekat
dengan sebuah garis lurus; korelasi tidak menunjukkan mana yang menyebabkan
yang mana, atau bahwa salah satunya menjadi penyebab. (b) Iklim yang lebih
hangat dapat saja melepaskan karbon dioksida ke udara; atau penyebab ketiga,
seperti perubahan orbit Bumi, dapat saja menggerakkan keduanya; atau
masing-masing dapat saja mendorong yang lain. (c) Waktu terjadinya perubahan,
yaitu mana yang lebih dahulu, dari penanggalan yang teliti atas es dan
udaranya; dan pemahaman tentang bagaimana karbon dioksida bekerja terhadap
radiasi Bumi (\cref{ch:g12:climate-change}).
\end{solution}

\begin{solution}{exo:g12:past-climates:12}
(a) $\num{1300000}/\num{41000} \approx 32$ siklus antara 2.6 dan 1.3 juta
tahun yang lalu; $\num{700000}/\num{98500} \approx 7$ siklus selama
\num{700000} tahun terakhir. (b) \num{41000} tahun: kemiringan sumbu; sekitar
\num{100000} tahun: bentuk orbit. (c) Satu siklus \num{41000} tahun disampel
sekitar $\num{41000}/\num{2000} \approx 20$ kali, cukup untuk menggambarkannya;
perubahan yang berlangsung seabad jatuh di antara dua sampel, atau di dalam
satu sampel, dan hampir tidak meninggalkan jejak.
\end{solution}

\begin{solution}{exo:g12:past-climates:13}
(a) $192/10 \approx 19$ dan $198/10 \approx 20$: sekitar 20 kali lebih cepat.
(b) Dibandingkan dengan \qty{4}{ppm} per abad: $192/4 = 48$ sampai $198/4 \approx 50$
kali; dibandingkan dengan \qty{42}{ppm}: $192/42 \approx 4.6$ sampai $198/42 \approx 4.7$
kali. Kenaikan sekarang sekitar 5 sampai 50 kali lebih cepat daripada kenaikan
PETM. (c) Kenaikannya $2000 - 900 = \qty{1100}{ppm}$: $1100/42 \approx 26$
abad, sekitar \num{2600} tahun dengan laju PETM yang tercepat;
$1100/195 \approx 5.6$ abad, kurang dari 600 tahun, dengan laju sekarang.
\end{solution}

\begin{solution}{exo:g12:past-climates:14}
(a) Sepuluh kali 1 sampai \qty{1.5}{\degreeCelsius}: 10 sampai
\qty{15}{\degreeCelsius} per seribu tahun, yaitu 0.10 sampai
\qty{0.15}{\degreeCelsius} per dekade. (b) Laju global sekarang,
\qty{0.20}{\degreeCelsius} per dekade, adalah 1.3 sampai 2 kali pemanasan
tercepat di atas Greenland. (c) Satu tempat dapat berubah jauh lebih besar
daripada rata-rata seluruh dunia, tempat perubahan di berbagai wilayah sebagian
saling meniadakan: laju global dibandingkan dengan laju global.
\end{solution}

\begin{solution}{exo:g12:past-climates:15}
(a) \qty{14}{m} dengan laju \qty{12}{m} per seribu tahun memerlukan $14/12 \approx
1.17$ ribu tahun, sekitar \num{1170} tahun: sekitar \num{9800} tahun yang lalu.
(b) Tidak: permukaan laut harus tetap setidaknya \qty{36}{m} di bawah
permukaan sekarang untuk menyatukan Kalimantan dengan Asia; begitu melewati
\qty{-36}{m} jembatan darat itu tenggelam, dan Kalimantan kembali menjadi
pulau. (c) $36/12 = 3$ ribu tahun, sekitar \num{3000} tahun; sebenarnya
kenaikan itu melambat setelah sekitar \num{8200} tahun yang lalu, sehingga
waktunya lebih lama.
\end{solution}

\begin{solution}{pb:g12:past-climates:1}
\textbf{1.} Es itu sendiri, melalui isotop beratnya, merekam suhu; udara dalam
gelembung-gelembungnya merekam karbon dioksida.

\textbf{2.} $5.5 - (-10.6) = \qty{16.1}{\degreeCelsius}$.

\textbf{3.} $299 - 174 = \qty{125}{ppm}$.

\textbf{4.} Titik-titiknya terletak dekat dengan sebuah garis lurus: kedua
rekaman itu naik dan turun bersama. Koefisien itu tidak mengatakan mana yang
menyebabkan yang mana.

\textbf{5.} Rekaman masa lalu digambar dengan masa tertua di salah satu sisi;
bila dibaca ke arah yang salah, pemanasan akan dikira pendinginan.

\textbf{6.} \omterm{def:g12:past-climates:glacial-cycle}{Periode glasial} adalah selang dingin dalam sebuah zaman es, ketika
lapisan es tumbuh jauh melampaui ukurannya sekarang; \omterm{def:g12:past-climates:glacial-cycle}{interglasial} adalah selang
hangat di antara dua \omterm{def:g12:past-climates:glacial-cycle}{periode glasial}, dengan lapisan es yang mendekati ukuran
sekarang atau lebih kecil.

\textbf{7.} Dalam sebuah \omterm{def:g12:past-climates:glacial-cycle}{interglasial}, yang dimulai sekitar \num{11000} tahun
yang lalu.

\textbf{8.} $\num{800000}/\num{98000} \approx 8$: sekitar delapan siklus.

\textbf{9.} Bentuk orbit.

\textbf{10.} Irama \num{41000} tahun, periode kemiringan sumbu.

\textbf{11.} $-0.4 - (-9.4) = \qty{9.0}{\degreeCelsius}$.

\textbf{12.} $264 - 190 = \qty{74}{ppm}$.

\textbf{13.} \qty{6}{\degreeCelsius}; perubahan di Antarktika, \qty{9.0}{\degreeCelsius},
lebih besar, sekitar 1.5 kali perubahan global.

\textbf{14.} $\num{18000} - \num{11000} = \num{7000}$ tahun.

\textbf{15.} $6/7 \approx \qty{0.86}{\degreeCelsius}$ per seribu tahun.

\textbf{16.} Seribu tahun adalah seratus dekade: $0.86/100 \approx
\qty{0.0086}{\degreeCelsius}$ per dekade.

\textbf{17.} $0.20 \times 100 = \qty{20}{\degreeCelsius}$ per seribu tahun.

\textbf{18.} $20/0.86 \approx 23$ kali.

\textbf{19.} $20/1.5 \approx 13$ dan $20/1 = 20$: 13 sampai 20 kali.

\textbf{20.} Pemanasan global sekarang \textbf{sekitar dua puluh kali lebih
cepat} daripada pemanasan keluar dari zaman es terakhir: 23 kali laju
rata-ratanya, 13 sampai 20 kali laju tercepatnya.
\end{solution}
